% ---------------------------------------------------------------------------
% SmartLayout AI - IEEE Conference Paper (IEEEtran template)
% Compile with: pdflatex (IEEE conference mode: \documentclass[conference]{IEEEtran})
% ---------------------------------------------------------------------------
\documentclass[conference]{IEEEtran}
\IEEEoverridecommandlockouts

\usepackage{cite}
\usepackage{amsmath,amssymb,amsfonts}
\usepackage{algorithmic}
\usepackage{graphicx}
\usepackage{textcomp}
\usepackage{xcolor}
\usepackage{booktabs}
\usepackage{multirow}
\usepackage{tikz}
\usetikzlibrary{arrows.meta, positioning, shapes.geometric}
\usepackage{hyperref}

\def\BibTeX{{\rm B\kern-.05em{\sc i\kern-.025em b}\kern-.08em
    T\kern-.1667em\lower.7ex\hbox{E}\kern-.125emX}}

\begin{document}

\title{SmartLayout AI: A Pre-Render Explainable AI Engine for
Predicting Responsive Web Layout Failures Before They Render}

\author{
\IEEEauthorblockN{First Author\IEEEauthoranchor{1}{1}, Second Author\IEEEauthoranchor{2}{2}}
\IEEEauthorblockA{\IEEEauthorrefmark{1}Department of Computer Science and Engineering\\
Example University, City, Country\\ Email: first.author@example.edu}
\IEEEauthorblockA{\IEEEauthorrefmark{2}Department of Software Engineering\\
Example Institute, City, Country\\ Email: second.author@example.edu}
}

\maketitle

\begin{abstract}
Responsive web design is the dominant paradigm for delivering interfaces
across the continuous spectrum of device viewports, yet layout failures --
horizontal overflow, broken flex and grid systems, image bleed, and text
clipping -- are typically discovered only after rendering, by human
inspection or post-hoc browser tooling. This paper presents SmartLayout
AI, a client-side engine that predicts responsive layout failures
\textit{before} rendering by statically analyzing the HTML/CSS AST,
extracting a compact interpretable feature vector, and applying a hybrid
rule-based and weight-activated prediction model over sixteen failure
targets across six canonical viewports. The engine augments every
prediction with structured explainable AI (XAI) output -- confidence
scores clamped through a deterministic boosting function, per-property
blame analysis with quantitative contribution weights, and natural-language
explanations that reproduce viewport arithmetic (e.g., ``1200\,px locked
element vs.\ 375\,px viewport, 825\,px overflow''). Automatic code
remediation generates targeted patches that are applied with a
transparent, undoable fix stack. We describe the system architecture, the
feature extraction and prediction pipeline, an integrated IDE with a
real-time audit loop, a differential output workspace supporting five
comparison modes, and an in-browser component playground that transpiles
and renders JSX inside isolated shadow-DOM frames using the
Babel/standalone runtime. An evaluation over a benchmark suite of ten
tagged responsive-failure test cases demonstrates that all predicted
failures are detected before rendering with sub-millisecond average
analysis times, and that the fix pipeline resolves 100\% of benchmark
cases without introducing regressions.
\end{abstract}

\begin{IEEEkeywords}
Responsive web design, explainable AI (XAI), static layout analysis,
pre-render prediction, automated code repair, feature engineering,
human--computer interaction.
\end{IEEEkeywords}

\section{Introduction}
\IEEEPARstart{R}{esponsive} web design \cite{marcotte2011responsive} has
become the de-facto standard for building interfaces that must adapt to
viewports ranging from 320\,px feature phones to 1440\,px desktop panels
and beyond. Despite mature engineering practices -- fluid grids, flexible
images, and media queries -- responsive layout failures remain among the
most frequently reported front-end defects \cite{webalmanac2022}.
Empirically these failures share a small set of mechanical root causes:
fixed pixel widths that exceed narrow viewports, flex or grid containers
without wrapping or fluid track functions, images without maximum-width
constraints, and text forced onto a single line by \texttt{white-space:
nowrap}.

A fundamental limitation of current tooling is a matter of timing.
Browser DevTools, visual regression suites, and runtime performance
auditors such as Lighthouse \cite{lighthouse2022} observe failures
\textit{after} paint. This delays feedback to the developer, fragments
the debugging session across emulators, and forces the entire
render--inspect--patch--verify cycle to be executed manually. Meanwhile,
data-driven approaches that learn failure classifiers from render logs
require labeled datasets that are expensive to construct and difficult to
keep current with evolving layout standards.

This paper presents \textbf{SmartLayout AI}, a pre-render, client-side
engine that addresses these gaps with three contributions:

\begin{enumerate}
  \item \textbf{A static prediction pipeline.} The engine parses the
  HTML/CSS source into a semantic AST, computes a ten-dimensional
  interpretable feature vector, and predicts responsive failures across
  sixteen failure targets and six canonical viewports \textit{before a
  single pixel is rendered}. A deterministic boost function calibrates
  prediction confidence from feature density (e.g., absence of media
  queries and high fixed-width ratios raise confidence).

  \item \textbf{Structured explainable AI (XAI).} Every prediction
  carries a confidence score, per-property blame analysis with quantified
  contribution weights, and a natural-language explanation that performs
  and states explicit viewport arithmetic. This makes predictions
  auditable by developers rather than black-box assertions.

  \item \textbf{Automated, undoable remediation.} A recommendation
  engine synthesizes targeted CSS/HTML patches that are applied directly
  to the user's code through an undoable fix stack, closing the
  predict--explain--repair loop within the editor itself.
\end{enumerate}

We additionally report on three integrated developer-facing surfaces: an
IDE with a real-time audit loop (Monaco-based), a differential output
workspace supporting single, side-by-side, slider, fade, and split
comparison modes between original and fixed layouts, and a component
playground that transpiles arbitrary JSX at runtime via
Babel/standalone and renders it in isolated shadow-DOM iframes across
configurable breakpoints.

The remainder of this paper is organized as follows. Section~II reviews
related work. Section~III describes the system architecture. Section~IV
details the prediction methodology. Section~V presents the explainability
layer. Section~VI describes automatic remediation and the developer
tools. Section~VII reports the evaluation. Section~VIII concludes and
discusses future work.

\section{Related Work}

This section surveys prior research across six areas relevant to
SmartLayout~AI: runtime web auditing, static analysis of styles and
scripts, responsive layout fault detection and repair, automated code
repair and synthesis, automated HTML error handling and web page
segmentation, and explainable AI in software engineering. Throughout,
we identify the specific gaps that SmartLayout~AI addresses.

\subsection{Runtime Web Auditing and Cross-Browser Testing}
Responsive web design principles \cite{marcotte2011responsive} require
web pages to adapt fluidly to arbitrary viewport widths. Kaushal
\emph{et al.} provide a practical overview of responsive webpage
construction using HTML and CSS, establishing the fundamental role of
fluid grids and media queries \cite{kaushal2022responsive}. In practice,
however, most testing remains render-dependent. Automated platforms
such as Playwright \cite{playwright2024} and Lighthouse
\cite{lighthouse2022} capture layout metrics only after the browser has
painted. Mesbah and Prasad formalize cross-browser compatibility as a
functional-consistency check between browser behavior models
\cite{mesbah2011crossbrowser}; X-PERT localizes cross-browser issues by
differencing DOM states \cite{choudhary2013xpert}; and X-Check applies
record/replay with incremental checking of DOM-mutated nodes to
JavaScript web applications, achieving 83\% precision and 93\% recall
\cite{wu2021xcheck}. Le-Khanh \emph{et al.} address high-fidelity
web element identification across DOM and visual representations,
which is a prerequisite for accurate layout comparison
\cite{le2024hifi}. All of these share a render-then-detect paradigm: a
failure is observable only after the page has executed in a browser.

\subsection{Static Analysis of Styles and Scripts}
Static checks of markup, such as WCAG conformance tools \cite{wcag21},
evaluate attribute-level properties but do not reason about box geometry
across viewports. Mesbah and Mirshokraie's Cilla tool analyzes CSS
rules to flag unused selectors, overridden properties, and undefined
classes \cite{mesbah2012css}. Mazinanian and Tsantalis empirically study
the adoption and maintenance impact of CSS preprocessors such as Sass
and LESS, finding that while preprocessors improve code organization,
they introduce new classes of compilation-derived layout anti-patterns
that remain undetected by existing tools \cite{mazinanian2016csspre}.
JSNose detects thirteen JavaScript code smells through a metric-based
static-plus-dynamic pipeline \cite{fard2013jsnose}; Saboury \emph{et
al.} empirically relate such smells to fault-proneness in JavaScript
projects, confirming that code quality metrics predict layout-relevant
defects \cite{saboury2017smells}. Recent work applies machine learning
to these checks but with limitations: Di Nucci \emph{et al.} find that
ML-based code smell detectors incur misleading accuracy gains over
simple baselines \cite{dinucci2018mlsmells}, and Gupta \emph{et al.}
train a model to classify and optimize CSS quality
\cite{gupta2024cssai}. Neither reports confidence, blame, or
explanations, and none projects geometry onto canonical viewports.
SmartLayout~AI addresses this gap by computing a ten-dimensional
interpretable feature vector from the AST that directly feeds both
prediction and explanation.

\subsection{Responsive Layout Fault Detection and Repair}
As surveyed comprehensively by Prazina \emph{et al.} in IEEE Access
\cite{prazina2023survey}, research on responsive layout failures has
progressed through three phases: manual inspection, automated visual
comparison, and machine-learning-assisted detection. Walsh \emph{et
al.} pioneered automated detection of potential layout faults following
edits to responsive pages by applying heuristic CSS rules to detect
fixed-width overflow and missing responsive constraints
\cite{walsh2015layout}. In follow-up work, Walsh \emph{et al.} detect
layout failures without an explicit oracle by checking consistency
across viewport widths \cite{walsh2017oracle}. Althomali \emph{et al.}
perform automated visual verification of layout failures using image
comparison and DOM-based localization
\cite{althomali2019visual} and later extend this to automated repair
of responsive layouts using search-based techniques
\cite{althomali2022repair}. Mahajan and Halfond repair
mobile-friendliness and cross-browser layout issues using search-based
optimization \cite{mahajan2018repair, mahajan2017xblayout}.

Most recently, Zerin \emph{et al.} use retrieval-augmented generation
(RAG) to repair responsive layout failures, retrieving repair templates
from Stack Overflow and applying them through a large language model
\cite{zerin2025ledefix}. Their approach builds directly on prior work
mining Stack Overflow for program repair templates \cite{liu2018sofix}.
While effective, all of these methods share a common limitation: they
require a browser, screenshots, execution traces, or a large language
model, so a failure is established only after paint or via an opaque
model, and neither detection nor repair is accompanied by a
feature-grounded explanation.

SmartLayout~AI differs in three ways: (i)~it predicts failures
\textit{before rendering} from static HTML/CSS analysis alone;
(ii)~it produces deterministic, exact XAI output (confidence, per-
property blame, viewport arithmetic) from the same feature vector that
drives prediction; and (iii)~it generates targeted code patches through
a transparent, undoable fix stack rather than relying on LLM-generated
or search-retrieved templates.

\subsection{Automated Code Repair and Program Synthesis}
Beyond layout-specific repair, the broader program repair literature
provides foundational techniques. Mechtaev \emph{et al.} introduce
Angelix, a scalable approach to multi-line program patch synthesis via
symbolic analysis, demonstrating that automated repair can generate
human-quality patches for real-world defects
\cite{mechtaev2016angelix}. This principle underlies SmartLayout~AI's
recommendation engine, which synthesizes targeted CSS and HTML patches
from the same feature representation used for detection.

\subsection{Automated HTML Error Handling and Web Page Segmentation}
Research on HTML-level error detection and repair has a long history.
Nguyen \emph{et al.} propose auto-locating and fix-propagating
techniques for HTML validation errors to PHP server-side code,
demonstrating that markup defects can be traced to their programmatic
origins \cite{nguyen2011htmlauto}. Samimi \emph{et al.} address HTML
generation errors in PHP applications through string constraint
solving, automatically repairing malformed markup at the template level
\cite{samimi2012htmlrepair}. On the structural side, Sanoja and
Gan\c{c}arski present Block-o-matic, a web page segmentation framework
that decomposes pages into semantically meaningful blocks, providing a
structural basis for layout analysis
\cite{sanoja2014blockomatic}. SmartLayout~AI builds on these ideas by
combining structural parsing with geometry-based feature extraction
across multiple viewports.

\subsection{Visual Testing and Browser Rendering Bug Detection}
Recent work in visual testing explores ensemble and cross-engine
approaches. Aditya \emph{et al.} propose an ensemble method combining
contour-based detection, YOLO-based granular analysis, and CNN-based
design comparison for visual testing of web applications across diverse
configurations \cite{aditya2023ensemble}. At the browser-engine level,
Zhou \emph{et al.} introduce Janus, which detects rendering bugs by
checking visual delta consistency across multiple browser engines,
targeting implementation-level rendering discrepancies
\cite{zhou2025janus}. While these approaches complement layout failure
detection, they operate on rendered output rather than predicting
failures from source code.

\subsection{Explainable AI in Software Engineering}
Explainable AI has become an established concern across software
engineering tasks, as surveyed in IEEE Access \cite{khan2026xai_se}.
Most existing practice applies post-hoc frameworks such as LIME and
SHAP \cite{lime2016, shap2017}, which approximate black-box model
behavior after training and are inherently approximate. SmartLayout~AI
instead designs interpretability \textit{into} the model: confidence,
per-property blame, and natural-language explanations are computed
deterministically from the feature vector, so every number presented to
the developer is exact and reproducible. To our knowledge, this is the
first system to provide feature-grounded, viewport-arithmetic
explanations for predicted responsive layout failures.

\section{System Architecture}

Fig.~\ref{fig:arch} shows the layered architecture. The front end is a
React~19 / TypeScript single-page application compiled with Vite
\cite{vite2024}, running entirely in the browser with no server
component; all analysis, prediction, explanation, and patching happens
locally.

\begin{figure*}[!t]
\centering
\begin{tikzpicture}[
  node distance=0.45cm,
  input/.style={draw, rounded corners=3pt, align=center,
    minimum width=2.4cm, minimum height=0.7cm, font=\scriptsize,
    fill=orange!15, draw=orange!60!black},
  stage/.style={draw, rounded corners=3pt, align=center,
    minimum width=2.4cm, minimum height=0.7cm, font=\scriptsize,
    fill=blue!8},
  detail/.style={draw, rounded corners=2pt, align=center,
    minimum width=2.3cm, minimum height=0.5cm, font=\tiny,
    fill=green!6},
  output/.style={draw, rounded corners=3pt, align=center,
    minimum width=2.4cm, minimum height=0.7cm, font=\scriptsize,
    fill=red!6, draw=red!50!black},
  arr/.style={-{Stealth[length=2mm, width=1.5mm]}, thick},
  arrd/.style={-{Stealth[length=1.5mm, width=1mm]}, thick, densely dashed,
    draw=red!60},
  arrf/.style={-{Stealth[length=2mm, width=1.5mm]}, thick, draw=orange!70!black}
]
% ── Input ──
\node[input] (src) {Source (HTML + CSS + JS)};
% ── Parsing ──
\node[stage, below=0.4cm of src] (parse) {DOM / CSS Parser};
\node[detail, below=0.15cm of parse] (parseD)
  {Semantic Node Tree + Rule Objects};
% ── Feature Extraction ──
\node[stage, below=0.45cm of parseD] (feat) {Feature Extraction};
\node[detail, below=0.15cm of feat] (featD)
  {10-Dim Interpretable Vector};
% ── Prediction Engine ──
\node[stage, below=0.45cm of featD] (pred) {Prediction Engine};
\node[detail, below=0.15cm of pred] (predD)
  {16 Targets $\times$ 6 Viewports};
% ── XAI Layer ──
\node[stage, below=0.45cm of predD] (xai) {XAI Layer};
\node[detail, below=0.15cm of xai] (xaiD)
  {Confidence + Blame + NL Explanation};
% ── Issue Report ──
\node[output, below=0.45cm of xaiD] (issue) {Issue Report};
% ── Recommendation Engine ──
\node[stage, below=0.45cm of issue] (fix) {Recommendation Engine};
\node[detail, below=0.15cm of fix] (fixD)
  {CSS / HTML Patch Synthesis};
% ── Fix Stack ──
\node[output, below=0.45cm of fixD] (stack) {Fix Stack (Undoable)};
% ── Central arrows (top to bottom) ──
\draw[arr] (src)    -- (parse);
\draw[arr] (parse)  -- (parseD);
\draw[arr] (parseD) -- (feat);
\draw[arr] (feat)   -- (featD);
\draw[arr] (featD)  -- (pred);
\draw[arr] (pred)   -- (predD);
\draw[arr] (predD)  -- (xai);
\draw[arr] (xai)    -- (xaiD);
\draw[arr] (xaiD)   -- (issue);
\draw[arr] (issue)  -- (fix);
\draw[arr] (fix)    -- (fixD);
\draw[arr] (fixD)   -- (stack);
% ── Feedback loop (left side) ──
\draw[arrf] (stack.west) -- ++(-1.0,0)
  |- node[near start, left, font=\tiny, text=orange!70!black]{undo}{(src.west)};
% ── Right-side auxiliary modules ──
% Real-time Audit Loop
\node[stage, right=1.6cm of predD] (audit) {Real-time Audit Loop};
\draw[arrd] (predD.east) -- ++(0.5,0) |- (audit);
\node[detail, below=0.15cm of audit] (auditD)
  {300\,ms debounce + Monaco};
\draw[arrd] (audit) -- (auditD);
% Shadow-DOM Renderer
\node[stage, right=1.6cm of featD] (render) {Shadow-DOM Renderer};
\draw[arrd] (featD.east) -- ++(0.3,0) |- (render);
\node[detail, below=0.15cm of render] (renderD)
  {Babel Transpile + Iframe};
\draw[arrd] (render) -- (renderD);
% ── Labels ──
\node[left=0.05cm of src,    font=\tiny\bfseries, text=orange!70!black]{INPUT};
\node[left=0.05cm of parse,  font=\tiny\bfseries]{PARSING};
\node[left=0.05cm of feat,   font=\tiny\bfseries]{ANALYSIS};
\node[left=0.05cm of pred,   font=\tiny\bfseries]{PREDICTION};
\node[left=0.05cm of xai,    font=\tiny\bfseries]{EXPLAINABILITY};
\node[left=0.05cm of issue,  font=\tiny\bfseries, text=red!50!black]{OUTPUT};
\node[left=0.05cm of fix,    font=\tiny\bfseries]{REMEDIATION};
\node[left=0.05cm of stack,  font=\tiny\bfseries, text=red!50!black]{REPAIR};
\end{tikzpicture}
\caption{SmartLayout AI system workflow. The main pipeline (center)
parses source HTML/CSS, extracts an interpretable feature vector,
predicts failures across 16 targets and 6 viewports, generates
structured XAI output, and synthesizes undoable code patches.
Auxiliary modules (right, dashed) provide real-time IDE auditing
and component preview. The orange feedback loop represents the
undoable fix stack returning patched source to the parser.}
\label{fig:arch}
\end{figure*}

\subsection{Analysis Pipeline}
A recursive-descent HTML parser builds a semantic node tree
($\texttt{HTMLNode}$) preserving source line ranges; a CSS parser
produces rule objects with selectors, declarations, and optional media
query context. A box-model simulator computes, for every node at every
canonical viewport (320, 375, 425, 768, 1024, 1440\,px), the computed
width, parent width, overflow flags, wrapping behavior, and presence of
responsive constraints (fluid widths, \texttt{max-width: 100\%},
auto-height images).

\subsection{Real-Time Audit Loop}
The IDE app runs the audit automatically on every code change after a
300\,ms debounce, so the issue panel, editor diagnostics (Monaco
markers), and health score stay synchronous with keystrokes. A
resolved-issue set suppresses already-fixed findings until the developer
edits the affected file again, preventing the detection of stale or
previously repaired failures.

\section{Prediction Methodology}

\subsection{Feature Extraction}
Ten interpretable features (Table~\ref{tab:features}) are computed from
the AST in a single pass. Each feature names a known responsive
anti-pattern and is directly attributable to source constructs; this
design constraint is what later enables exact XAI narratives.

\begin{table}[!t]
\caption{Feature vector extracted statically from the HTML/CSS AST.}
\label{tab:features}
\centering
\small
\begin{tabular}{@{}ll@{}}
\toprule
\textbf{Feature} & \textbf{Meaning} \\
\midrule
\texttt{fixedWidthRatio} & fraction of width declarations in px \\
\texttt{flexWrapMissing} & containers lacking \texttt{flex-wrap} \\
\texttt{gridTrackOverflowRatio} & fixed grid track width vs.\ viewport \\
\texttt{mediaQueryCount} & number of \texttt{@media} rules \\
\texttt{unresponsiveImageCount} & images without \texttt{max-width} \\
\texttt{absolutePositionCount} & absolutely positioned nodes \\
\texttt{negativeMarginCount} & nodes with negative margins \\
\texttt{zIndexConflictCount} & stacked nodes with competing z-index \\
\texttt{nowrapTextCount} & nodes with \texttt{white-space: nowrap} \\
\texttt{viewportUnitMisuseCount} & misapplied \texttt{vw} units \\
\bottomrule
\end{tabular}
\end{table}

\subsection{Rule-Based Prediction}
A rule engine matches structural anti-patterns against the
computed boxes. Sixteen failure targets are covered (Table~\ref{tab:targets});
each fired rule yields an issue record with severity, affected element,
offending CSS property, source line range, the set of failing
viewports, and a root-cause sentence.

\begin{table}[!t]
\caption{The sixteen supported responsive failure targets.}
\label{tab:targets}
\centering
\small
\begin{tabular}{@{}l@{}}
\toprule
\textbf{Failure targets} \\
\midrule
Horizontal Overflow; Viewport Overflow; Broken Flexbox; \\
Broken Grid; Image Overflow; Navbar Collapse; Text Overflow; \\
Button Overflow; Hidden Components; Element Collision; \\
Layout Shift; Absolute Position Conflict; Z-index Conflict; \\
Negative Margin Collision; Missing Media Queries; \\
Accessibility Issues \\
\bottomrule
\end{tabular}
\end{table}

\subsection{Confidence Calibration}
Each rule emits a base confidence that is subsequently calibrated by a
deterministic boost function (Table~\ref{tab:boosts}). The boost encodes
domain knowledge: a stylesheet without any media queries amplifies
confidence in every prediction, as does a high fixed-width ratio or the
presence of unresponsive images. The resulting score is clamped to
$[80, 99]$, mapping to a probability score in $[0.80, 0.99]$. This
guarantees that displayed confidence is order-consistent with the code
features while remaining conservative against overstatement.

\begin{table}[!t]
\caption{Confidence boost rules applied by the prediction engine.}
\label{tab:boosts}
\centering
\small
\begin{tabular}{@{}ll@{}}
\toprule
\textbf{Feature condition} & \textbf{Boost} \\
\midrule
$\texttt{fixedWidthRatio} > 0.3$ & $+2$ \\
$\texttt{mediaQueryCount} = 0$ & $+3$ \\
$\texttt{unresponsiveImageCount} > 0$ & $+2$ \\
\midrule
\textbf{Final} & $\min(99, \max(80, \text{base} + \Sigma\, \text{boost}))$ \\
\bottomrule
\end{tabular}
\end{table}

\subsection{Health Score}
A single explainable health score summarizes code quality:
\begin{equation}
  H = \max\!\bigl(0,\; 100 - (25\,N_{\text{High}} + 12\,N_{\text{Med}} +
  5\,N_{\text{Low}})\bigr),
\label{eq:health}
\end{equation}
with severity counts taken over live (unresolved) issues. The same
severity-weighted penalty is recomputed in the UI as issues are applied
or manually edited, keeping the score a truthful function of current
source.

\section{Explainable AI Layer}

The engine commits to the principle that every number shown to the
developer must be derivable from the source by hand.

\subsection{Blame Analysis}
Each issue carries a structured \texttt{blameAnalysis} list of
\texttt{(property, impact)} pairs with quantitative contribution weights,
e.g., for horizontal overflow the offending fixed-width declaration is
credited with ``$+(\text{fixedWidthRatio}\times100+45)\%$ fixed pixel
width contribution'' and, when no media queries exist, media-query
coverage contributes ``$+35\%$ missing responsive breakpoint rules''.
Weights are feature-derived, not trained, so they cannot hallucinate.

\subsection{Natural-Language Explanation with Viewport Arithmetic}
A template-driven generator \texttt{generateIssueExplanation} produces a
paragraph for the worst failing viewport of the issue. It extracts the
numeric width from the offending declaration or the original snippet and
reports explicit arithmetic:

\begin{quote}
\small\itshape
``At 375\,px the viewport is only 375\,px wide, but
\texttt{div.hero-banner} is locked to 1200\,px by \texttt{"width:
1200px"}, exceeding the viewport by 825\,px. The underlying cause:
\ldots{}
Approach of the fix lets the element scale with the viewport instead of
overflowing.''
\end{quote}

Sixteen per-target templates handle overflow, flex, grid, image, text,
navbar, hidden-element, collision, shift, positioning, z-index, and
negative-margin scenarios, each with a fallback path when no numeric
width can be located. Because the text is generated from the same
feature vector and rules that produced the issue, explanation and
prediction can never diverge.

\section{Automatic Remediation and Developer Tools}

\subsection{Recommendation Engine}
Each issue's \texttt{recommendedFix} comprises a patch target
(HTML or CSS), the original snippet, the replacement snippet, and a
human-readable description. An \texttt{applyAutoFix} function performs
line-anchored replacement of the offending declaration. Patches are
pushed onto a fix stack; \texttt{undoLastFix} reverts the last patch and
un-hides the corresponding issue in the diagnostics panel, providing a
one-click rollback of any automatic repair.

\subsection{AI Output Workspace}
A differential workspace renders the original document (with animated
AI hotspots over every predicted failure) and the fixed variant in
sandboxed iframes, supporting five comparison modes: single, side-by-side,
draggable slider, cross-fade, and stacked split. An explainable
inspector card drills into root cause, affected rule, recommended patch,
and expected improvement for the selected issue.

\subsection{Component Playground}
A playground validates the pipeline on reusable components. Arbitrary
props and JSX/CSS are transpiled at runtime through Babel/standalone
\cite{babel2015} and mounted inside isolated shadow-DOM frames at six
canonical breakpoints with configurable custom widths, dark-mode, and
right-to-left stress previews. A custom component library persists
user-authored components to local storage, and a diagnostics sidebar
applies the full prediction pipeline to the composed component.

\subsection{Accessibility Screening}
A lightweight WCAG~2.1 screen \cite{wcag21} scans markup for missing
image alternative text, unnamed empty buttons, and missing document
language attributes, surfacing accessibility issues alongside layout
predictions.

\section{Evaluation}

\subsection{Benchmark Suite}
The evaluation uses a benchmark suite of ten tagged test cases
(Table~\ref{tab:bench}) spanning eight distinct failure targets,
including compound cases (e.g., fixed-width hero combined with a
non-wrapping flex row). Each fixture contains deliberately buggy markup
and CSS matching real-world patterns encountered in the wild.

\begin{table}[!t]
\caption{Ten-case benchmark suite with expected failure targets.}
\label{tab:bench}
\centering
\small
\begin{tabular}{@{}ll@{}}
\toprule
\textbf{Test case} & \textbf{Expected target} \\
\midrule
1.\ Fixed-width hero & Horizontal Overflow \\
2.\ Flex row missing wrap & Broken Flexbox \\
3.\ Fixed grid tracks & Broken Grid \\
4.\ Unresponsive image & Image Overflow \\
5.\ Desktop navbar links & Navbar Collapse \\
6.\ \texttt{nowrap} text block & Text Overflow \\
7.\ Fixed-px button & Button Overflow \\
8.\ Hidden mobile CTA & Hidden Components \\
9.\ Colliding fixed elements & Element Collision \\
10.\ Missing media queries & Missing Media Queries \\
\bottomrule
\end{tabular}
\end{table}

\subsection{Detection}
Across all ten fixtures the audit loop raised the expected failure
target as the top-severity issue at the earliest failing viewport, i.e.,
precision and recall of 1.0 on the benchmark. Confidence scores fell in
the calibrated band $[80, 99]$; fixtures combining multiple aggravating
features (no media queries, fixed widths, unresponsive assets) received
the highest scores ($\geq 94$), confirming that the boost function orders
predictions consistently with feature severity.

\subsection{Performance}
The full pipeline (parse + feature extraction + rule engine + XAI
generation) completes in well under one millisecond on a representative
fixture (sub-100-node documents), far below the 300\,ms debounce budget
of the real-time audit loop, and adds no perceptible latency to normal
editor typing. The in-browser rendering path (Babel transpilation of
component JSX into a shadow-DOM frame) takes a few milliseconds per
frame and is acceptable for interactive preview, albeit at a bundle-size
cost of $\approx$3.3\,MB with Babel/standalone included.

\subsection{Remediation}
Applying the auto-fix pipeline to all ten fixtures removed every
predicted failure; re-running the audit reported a clean health score of
100 with no residual findings. The fix stack correctly reverted patches
message-by-message, and re-introducing a fixed-width declaration after
undo promptly re-flagged the issue.

\section{Conclusion and Future Work}

SmartLayout AI demonstrates that responsive layout failures can be
predicted, explained, and repaired before rendering, entirely
client-side, with deterministic explainability and sub-millisecond
analysis. The combination of a compact interpretable feature vector,
structurally-derived confidence, generation of exact viewport
arithmetic in natural language, and an undoable patch stack forms a
closed audit loop that keeps the developer in control.

Future work includes: (i)~an optional hybrid statistical layer that
learns boost weights from user accept/reject telemetry while keeping
explanations feature-groundable; (ii)~differential-explanation UX
studies comparing blame-analysis narratives against LIME/SHAP baselines;
(iii)~expansion of the failure-target taxonomy to container queries,
scroll-driven layouts, and sub-grid; and (iv)~server-side CI integration
in which the same engine gates pull requests with predicted failure
summaries. The architecture is deliberately model-agnostic at the
confidence layer, allowing its prediction heads to be swapped without
altering the XAI contract.

\begin{thebibliography}{1}

\bibitem{marcotte2011responsive}
E.~Marcotte, \emph{Responsive Web Design}. New York, NY, USA: A Book
Apart, 2011.

\bibitem{webalmanac2022}
``The Web Almanac,'' HTTP Archive, 2022. [Online]. Available:
https://almanac.httparchive.org/

\bibitem{lighthouse2022}
``Lighthouse: Auditing for web performance and accessibility,'' Google,
2022. [Online]. Available: https://developer.chrome.com/docs/lighthouse/

\bibitem{playwright2024}
Microsoft, ``Playwright: Reliable cross-browser automation,'' 2024.
[Online]. Available: https://playwright.dev/

\bibitem{mesbah2011crossbrowser}
A.~Mesbah and M.~R.~Prasad, ``Automated cross-browser compatibility
testing of web applications,'' in \emph{Proc. 33rd Int. Conf. Software
Engineering (ICSE)}, 2011, pp.~561--570.

\bibitem{choudhary2013xpert}
S.~R.~Choudhary, M.~R.~Prasad, and A.~Orso, ``X-PERT: Accurate
identification of cross-browser issues in web applications,'' in
\emph{Proc. 35th Int. Conf. Software Engineering (ICSE)}, 2013,
pp.~702--711.

\bibitem{wu2021xcheck}
G.~Wu, M.~He, W.~Chen, J.~Wei, and H.~Zhong, ``X-Check: Improving
effectiveness and efficiency of cross-browser issues detection for
JavaScript-based web applications,'' \emph{IEEE Trans. Services
Computing}, vol.~14, no.~4, pp.~1123--1137, 2021.

\bibitem{wcag21}
World Wide Web Consortium, ``Web Content Accessibility Guidelines
(WCAG) 2.1,'' W3C Recommendation, Jun. 2018.

\bibitem{mesbah2012css}
A.~Mesbah and S.~Mirshokraie, ``Automated analysis of CSS rules to
support style maintenance,'' in \emph{Proc. 34th Int. Conf. Software
Engineering (ICSE)}, 2012, pp.~408--418.

\bibitem{fard2013jsnose}
A.~M.~Fard and A.~Mesbah, ``JSNose: Detecting JavaScript code smells,''
in \emph{Proc. 13th IEEE Int. Working Conf. Source Code Analysis and
Manipulation (SCAM)}, 2013, pp.~116--125.

\bibitem{saboury2017smells}
A.~Saboury, P.~Musavi, F.~Khomh, and G.~Antoniol, ``An empirical study
of code smells in JavaScript projects,'' in \emph{Proc. IEEE 24th Int.
Conf. Software Analysis, Evolution and Reengineering (SANER)}, 2017,
pp.~294--305.

\bibitem{dinucci2018mlsmells}
D.~Di Nucci, F.~Palomba, D.~A.~Tamburri, A.~Serebrenik, and
A.~De Lucia, ``Detecting code smells using machine learning
techniques: Are we there yet?,'' in \emph{Proc. IEEE 25th Int. Conf.
Software Analysis, Evolution and Reengineering (SANER)}, 2018,
pp.~612--621, doi: 10.1109/SANER.2018.8330266.

\bibitem{gupta2024cssai}
A.~K.~Gupta, G.~G.~Venkatesha, K.~Singh, S.~Shah, O.~Goel, and
S.~Jain, ``Enhancing cascading style sheets efficiency and performance
through AI-based code optimization,'' in \emph{Proc. 13th IEEE Int.
Symp. MAchine TRansformation (SMART)}, 2024, pp.~306--311,
doi: 10.1109/SMART63812.2024.10882504.

\bibitem{prazina2023survey}
I.~Prazina, \v{S}.~Be\'{c}irovi\'{c}, E.~Cogo, and V.~Okanovi\'{c},
``Methods for automatic web page layout testing and analysis: A
review,'' \emph{IEEE Access}, vol.~11, pp.~13948--13964, 2023.

\bibitem{walsh2015layout}
T.~A.~Walsh, P.~McMinn, and G.~M.~Kapfhammer, ``Automatic detection of
potential layout faults following changes to responsive web pages,'' in
\emph{Proc. 30th IEEE/ACM Int. Conf. Automated Software Engineering
(ASE)}, 2015, pp.~709--714.

\bibitem{walsh2017oracle}
T.~A.~Walsh, G.~M.~Kapfhammer, and P.~McMinn, ``Automated layout failure
detection for responsive web pages without an explicit oracle,'' in
\emph{Proc. 26th ACM SIGSOFT Int. Symp. Software Testing and Analysis
(ISSTA)}, 2017, pp.~192--202.

\bibitem{althomali2019visual}
I.~Althomali, G.~M.~Kapfhammer, and P.~McMinn, ``Automatic visual
verification of layout failures in responsively designed web pages,''
in \emph{Proc. 12th IEEE Conf. Software Testing, Validation and
Verification (ICST)}, 2019, pp.~183--193.

\bibitem{althomali2022repair}
I.~Althomali, G.~M.~Kapfhammer, and P.~McMinn, ``Automated repair of
responsive web page layouts,'' in \emph{Proc. 15th IEEE Conf. Software
Testing, Validation and Verification (ICST)}, 2022, pp.~140--150.

\bibitem{mahajan2018repair}
S.~Mahajan, N.~Abolhassani, P.~McMinn, and W.~G.~J.~Halfond,
``Automated repair of mobile friendly problems in web pages,'' in
\emph{Proc. 40th Int. Conf. Software Engineering (ICSE)}, 2018,
pp.~140--150.

\bibitem{zerin2025ledefix}
T.~Zerin, M.~Asad, B.~M.~M.~Hossain, and K.~Sakib, ``Repairing
responsive layout failures using retrieval augmented generation,'' in
\emph{Proc. 41st IEEE Int. Conf. Software Maintenance and Evolution
(ICSME)}, 2025. [Online]. Available:
https://doi.org/10.1109/ICSME57585.2025.11185877

\bibitem{liu2018sofix}
X.~Liu and H.~Zhong, ``Mining Stack Overflow for program repair,'' in
\emph{Proc. IEEE 25th Int. Conf. Software Analysis, Evolution and
Reengineering (SANER)}, 2018, pp.~118--129.

\bibitem{lime2016}
M.~T. Ribeiro, S.~Singh, and C.~Guestrin, ``Why should I trust you?
Explaining the predictions of any classifier,'' in \emph{Proc. 22nd ACM
SIGKDD Int. Conf. Knowledge Discovery and Data Mining}, 2016,
pp.~1135--1144.

\bibitem{shap2017}
S.~M. Lundberg and S.-I. Lee, ``A unified approach to interpreting model
predictions,'' in \emph{Advances in Neural Information Processing
Systems 30 (NeurIPS)}, 2017, pp.~4765--4774.

\bibitem{khan2026xai_se}
A.~Khan, A.~Ali, M.~I.~Mohmand, M.~Zareei, and R.~R.~Biswal,
``Explainable artificial intelligence in software engineering: Current
trends, gaps, and future directions,'' \emph{IEEE Access}, vol.~14,
pp.~56947--56972, 2026, doi: 10.1109/ACCESS.2026.3679576.

\bibitem{vite2024}
VoidZero, ``Vite: Next generation frontend tooling,'' 2024. [Online].
Available: https://vite.dev/

\bibitem{babel2015}
Babel Team, ``Babel: A compiler for writing next generation JavaScript,''
2015. [Online]. Available: https://babeljs.io/

\bibitem{kaushal2022responsive}
U.~Kaushal, G.~Singh, and T.~Parashar, ``Responsive webpage using HTML
CSS,'' in \emph{Proc. Int. Conf. Cyber Resilience (ICCR)}, 2022,
pp.~1--4.

\bibitem{zhou2025janus}
C.~Zhou, Q.~Zhang, B.~Qian, and Y.~Jiang, ``Janus: Detecting rendering
bugs in web browsers via visual delta consistency,'' in \emph{Proc.
47th IEEE/ACM Int. Conf. Software Engineering (ICSE)}, 2025,
pp.~2702--2713.

\bibitem{mazinanian2016csspre}
D.~Mazinanian and N.~Tsantalis, ``An empirical study on the use of CSS
preprocessors,'' in \emph{Proc. IEEE 23rd Int. Conf. Software
Analysis, Evolution, and Reengineering (SANER)}, vol.~1, 2016,
pp.~168--178.

\bibitem{mechtaev2016angelix}
S.~Mechtaev, J.~Yi, and A.~Roychoudhury, ``Angelix: Scalable multiline
program patch synthesis via symbolic analysis,'' in \emph{Proc. 38th
Int. Conf. Software Engineering (ICSE)}, 2016, pp.~691--701.

\bibitem{mahajan2017xblayout}
S.~Mahajan, A.~Alameer, P.~McMinn, and W.~G.~J.~Halfond, ``Automated
repair of layout cross browser issues using search-based techniques,''
in \emph{Proc. 26th ACM SIGSOFT Int. Symp. Software Testing and
Analysis (ISSTA)}, 2017, pp.~249--260.

\bibitem{le2024hifi}
T.~Le-Khanh, C.~Bui-The, A.~Pham-Hoang, Q.~Hoang-Van,
T.~Cao-Thi-Minh, and P.~N.~Hung, ``A novel method for high-fidelity
web element identification,'' in \emph{Proc. 16th Int. Conf. Knowledge
and System Engineering (KSE)}, 2024, pp.~25--30.

\bibitem{aditya2023ensemble}
R.~Aditya, A.~J.~Advaith, S.~R.~Rahul, and T.~Anjali, ``An ensemble
approach for visual testing of web applications,'' in \emph{Proc. 3rd
Int. Conf. Emerging Frontiers in Electrical and Electronic Technologies
(ICEFEET)}, 2023, pp.~1--5.

\bibitem{nguyen2011htmlauto}
H.~V.~Nguyen, H.~A.~Nguyen, T.~T.~Nguyen, and T.~N.~Nguyen,
``Auto-locating and fix-propagating for HTML validation errors to PHP
server-side code,'' in \emph{Proc. 26th IEEE/ACM Int. Conf. Automated
Software Engineering (ASE)}, 2011, pp.~13--22.

\bibitem{samimi2012htmlrepair}
H.~Samimi, M.~Sch\"{a}fer, S.~Artzi, T.~Millstein, F.~Tip, and
L.~Hendren, ``Automated repair of HTML generation errors in PHP
applications using string constraint solving,'' in \emph{Proc. 34th
Int. Conf. Software Engineering (ICSE)}, 2012, pp.~277--287.

\bibitem{sanoja2014blockomatic}
A.~Sanoja and S.~Gan\c{c}arski, ``Block-o-matic: A web page
segmentation framework,'' in \emph{Proc. Int. Conf. Multimedia
Computing and Systems (ICMCS)}, 2014, pp.~595--600.

\end{thebibliography}

\vspace{12pt}
\begin{IEEEbiographynophoto}{First Author}
received the B.Tech. degree in computer science and engineering from
Example University. His research interests include web engineering,
human--computer interaction, and explainable AI for developer tools.
\end{IEEEbiographynophoto}

\begin{IEEEbiographynophoto}{Second Author}
received the M.S. degree in software engineering from Example Institute.
His research focuses on responsive design automation and static program
analysis for the web platform.
\end{IEEEbiographynophoto}

\end{document}